\chapter{أدوات البناء}
\label{ch:builder}

\section{ماذا ستنجز؟}

تتيح لك أدوات البناء أن تصمّم تطبيقاً خاصاً بعملك عندما لا تجد في التطبيقات الجاهزة سجلاً يناسبه، مثل طلب داخلي أو سجل معدات أو متابعة زيارات. تبدأ بتحديد ما يمثله السجل الواحد، ثم تضيف حقوله وترتّب صفحته وقوائمه، وبعد ذلك تنشره ليستخدمه الفريق. ويمكنك لاحقاً أن تضيف إليه نموذج طلب، ومسار اعتماد، ومستندات تُنشأ من بياناته، ولوحة أرقام تتابعه.

يتولى هذا العمل مسؤول النظام أو من يملك صلاحية البناء، ويحتاج إلى مقعد كامل؛ إذ لا تظهر أدوات البناء لصاحب مقعد الموظف. أما الموظفون فيستخدمون التطبيق بعد نشره دون أن يحتاجوا إلى أدوات البناء.

\begin{rcnote}[مثال: شركة الأفق]
تريد شركة الأفق أن تستقبل طلبات زيارة الفروع في مكان واحد بدلاً من الرسائل المتفرقة. لذلك يبني مسؤول النظام نوع سجلات باسم «طلبات الزيارة»، يرفع فيه عمر حداد طلب زيارة فرع حلب، وتراجعه سارة الخطيب.
\end{rcnote}

\begin{figure}[H]
\centering
\begin{tikzpicture}[node distance=0.5cm]
\node[stepBox] (a) {\RL{تحديد ما يمثله السجل}};
\node[stepBox,below=of a] (b) {\RL{إضافة الحقول}};
\node[stepBox,below=of b] (c) {\RL{ترتيب الصفحة والقوائم}};
\node[stepBox,below=of c] (d) {\RL{القواعد والصلاحيات}};
\node[stepBox,below=of d] (e) {\RL{المراجعة والنشر}};
\node[stepBox,below=of e] (f) {\RL{إظهاره في التنقل}};
\node[stepBox,below=of f] (g) {\RL{نماذج ومستندات وتقارير}};
\draw[flowsTo] (a) -- (b);
\draw[flowsTo] (b) -- (c);
\draw[flowsTo] (c) -- (d);
\draw[flowsTo] (d) -- (e);
\draw[flowsTo] (e) -- (f);
\draw[flowsTo] (f) -- (g);
\end{tikzpicture}
\rclegend
\caption{مسار بناء تطبيق جديد}
\label{fig:builder-overview-flow}
\shotnote{يبيّن الشكل الترتيب الذي يوصي به هذا الفصل: تبني نوع السجلات أولاً، لأن كل ما بعده يعتمد على حقوله. ثم تنشره وتضعه في التنقل ليصل إليه الفريق، وبعد أن يعمل بنجاح تضيف النماذج والمستندات والتقارير. ولا يرى الفريق أي تعديل قبل النشر، لذلك جرّب كل مرحلة بسجل فعلي قبل الانتقال إلى التالية.}
\end{figure}

\section{كيف تترابط أدوات البناء؟}
\label{sec:builder-map}

قد تبدو أدوات البناء كثيرة، لكنها تدور حول فكرة واحدة: نوع السجلات. فهو المكان الذي تُحفظ فيه البيانات، وكل أداة أخرى تتعامل معه بطريقة مختلفة: النموذج يجمع البيانات إليه، والإجراءات تشغّل العمل عليه، والقوالب تُخرج منه مستندات، ولوحات المعلومات تعرض أرقامه.

\begin{figure}[H]
\centering
\begin{tikzpicture}[node distance=0.5cm]
\node[entity] (type) {\RL{نوع السجلات}};
\node[entityOpt,below right=0.8cm and 0.6cm of type] (forms) {\RL{النماذج}\\\RL{تجمع البيانات}};
\node[entityOpt,below=0.8cm of type] (sop) {\RL{الإجراءات}\\\RL{تشغّل العمل}};
\node[entityOpt,below left=0.8cm and 0.6cm of type] (out) {\RL{القوالب واللوحات}\\\RL{تُخرج النتائج}};
\draw[dependsOnOpt] (type.south) -- ++(0,-0.35) -| (forms.north);
\draw[dependsOnOpt] (type.south) -- (sop.north);
\draw[dependsOnOpt] (type.south) -- ++(0,-0.35) -| (out.north);
\end{tikzpicture}
\rclegend
\caption{الصورة الكاملة: كل شيء يدور حول نوع السجلات}
\label{fig:builder-big-picture}
\shotnote{يضع الشكل نوع السجلات في الأعلى، لأن بقية الأدوات تحتاج إليه، وكلها اختيارية. فالنماذج تُدخل البيانات، والإجراءات تشغّل العمل عليها، والقوالب واللوحات تُخرج منها مستندات وأرقاماً. فقد يكفيك نوع سجلات يُدخل الفريق بياناته مباشرة، ثم تضيف نموذجاً أو قالباً أو لوحة عندما تظهر الحاجة. لذلك لا تبدأ بالأتمتة أو التقارير قبل أن يعمل نوع السجلات بسجلات حقيقية.}
\end{figure}

\begin{table}[H]
\centering
\caption{كلمات ستقابلها في هذا الفصل}
\label{tab:builder-terms}
\begin{tabularx}{\linewidth}{R{3.6cm} Y}
\rcHeaderRow
\thh{الكلمة} & \thh{معناها ببساطة}\\
نوع السجلات (الكيان) & قالب لمجموعة سجلات متشابهة، مثل «طلبات الزيارة». وكل طلب منها سجل.\\
الحقل & معلومة واحدة في السجل، مثل تاريخ الزيارة.\\
المنشئ & المكان الذي تعدّل فيه نوع السجلات، وأقسامه هي الحقول وصفحة السجل والقوائم والقواعد والصلاحيات وغيرها.\\
المسودة والنشر & يُحفظ كل تعديل مسودةً لا يراها أحد، ثم يصل إلى الفريق عند النشر في إصدار جديد.\\
الدور والفعل & الدور مجموعة صلاحيات تُمنح للأشخاص، والفعل عمل مسموح به، مثل الإضافة أو التغيير.\\
النموذج & صفحة أسئلة يعبّئها الموظف أو مقدّم الطلب، وقد تكتب إجاباتها في نوع سجلات.\\
الإجراء الداخلي & مسار عمل من خطوات، مثل طلب اعتماد أو إنشاء مهمة، يبدأ بحدث تحدده.\\
القالب & ملف Word أو Excel فيه وسوم تُستبدل ببيانات السجل لإنشاء مستند.\\
\end{tabularx}
\end{table}

\section{دليل أدوات البناء}
\label{sec:builder-pages}

أدوات البناء واسعة، وفي كل أداة منها تفاصيل كثيرة؛ لذلك أفردنا لها دليلاً مستقلاً تبدأ فصوله من 1. يشرح كل فصل أداةً واحدة من أولها إلى آخرها، ويقسّمها إلى صفحات قصيرة، في كل صفحة مثال من شركة الأفق وصور مرقمة وخطوات للتحقق. ابدأ بفصل «الكيانات»، لأن الأدوات الأخرى تعمل على أنواع السجلات التي تبنيها فيه، ثم انتقل إلى بقية الفصول بحسب حاجة عملك.

\begin{rcnote}[مثال: شركة الأفق]
يبدأ مسؤول النظام في الأفق بفصل «الكيانات» ليبني «طلبات الزيارة» وينشرها، ثم ينتقل إلى فصل «التطبيقات والتنقل» ليضعها في تطبيق «خدمات الفروع». وعندما يحتاج عمر حداد إلى نموذج يرفع منه الطلب، يفتح المسؤول فصل «النماذج».
\end{rcnote}

\begin{table}[H]
\centering
\caption{فصول دليل أدوات البناء}
\label{tab:guide-builder-chapters}
\rcZebra
\begin{tabularx}{\linewidth}{R{4.2cm} Y}
\rcHeaderRow
\thh{الفصل} & \thh{ماذا يشرح؟}\\
الكيانات & إنشاء نوع السجلات وبناؤه في «المنشئ»: الحقول وصفحة السجل والقوائم والقواعد والصلاحيات، ثم مراجعته ونشره.\\
التطبيقات والتنقل & ترتيب التطبيقات والأقسام التي يراها الأعضاء، وما يظهر لكل عضو أولاً.\\
النماذج & بناء نماذج يعبّئها الموظفون أو مقدّمو الطلبات، وربط إجاباتها بأنواع السجلات، ثم نشرها.\\
الخدمات & تقديم النماذج خدماتٍ يطلبها الجمهور أو الموظفون، ومتابعة الطلبات وأدائها.\\
الإجراءات الداخلية & بناء مسارات عمل من خطوات تبدأ بحدث تحدده، مثل طلب اعتماد أو إنشاء مهمة.\\
إجراءات التشغيل القياسية والبوتات & أتمتة تربط كيان بخدمات وحسابات أخرى، وبوتات تعمل في الدردشة والبريد والتقويم.\\
المستندات والطباعة & قوالب Word وExcel والطباعة، والهويات والترويسات التي تظهر في المستندات.\\
لوحات المعلومات والتقارير & لوحات الأرقام والصفحات المستقلة، والتحليلات والتقارير بالذكاء الاصطناعي.\\
\end{tabularx}
\end{table}

أما «قوالب الحالات» و«إعدادات المهام» فتظهران أيضاً في قسم «الأساسية» من أدوات البناء، وقد شرحهما فصل المهام: الأولى في الصفحة \ref{sec:tasks-feature-statuses}، والثانية في الصفحة \ref{sec:tasks-feature-settings}.
